% --- NEURO PROBABILISTIC MODEL ---
\section{Neuro Probabilistic Model}

% --- SLIDE 1: MOTIVATION (Provided by you) ---
\begin{frame}{Neuro Probabilistic Model: Motivation}
    \begin{block}{Observation}
        Standard architectures provide \textbf{deterministic} predictions, lacking a \textbf{stochastic} explanation of confidence.
    \end{block}

    \begin{itemize}
        \item \textbf{Goal}: Model not only acceleration but also it's variance. 
        \item \textbf{Confidence Prediction}:
        \begin{itemize}
            \item Prediction is associated with a \textbf{confidence contour}.
            \item A \textbf{$95\%$ confidence region} around the expected value computed using the interval: $\ddot{q}_i \in [-2\sigma_i, 2\sigma_i]$
        \end{itemize}
        \item \textbf{Applications}: 
        \begin{itemize}
            \item \textbf{Diagnosis}: Enables anomaly detection (alarm triggers) when observations fall outside the predicted variance.
            \item \textbf{Trajectory generation}: Generate trajectory with an input and evaluate the confidence
            \item \textbf{Model Predictive Control}: Used to predict the next sliding window with safety margins
        \end{itemize}
        
    \end{itemize}
\end{frame}

% --- SLIDE 2: ARCHITECTURE (Based on ProbabilisticBaselineFNN) ---
\begin{frame}{Neuro Probabilistic Model: Architecture}
    To capture aleatoric uncertainty, we modify the Feed-Forward Network (FNN) to output parameters of a probability distribution rather than a single value.


    \begin{itemize}
        \item \textbf{Mean} ($\mu$): The expected acceleration.
        \item \textbf{Log-Variance} ($\log \sigma^2$): The estimated uncertainty. Note that we use log for stability. 
    \end{itemize}


    \begin{figure}[T]
            \centering
            \includegraphics[width=0.8\linewidth]{09_Neuro_Prob/neuro_prob_pred.pdf} \\
            \vspace{-1em}
    \end{figure}
    
\end{frame}


% --- SLIDE 3: MATHEMATICAL FORMULATION ---
\begin{frame}{Neuro Probabilistic Model: Probabilistic Formulation}
    The model assumes the ground truth acceleration $\ddot{q}$ follows a Gaussian distribution conditioned on the state:
    
    $$ p(\ddot{q} | \mathbf{q}) = \mathcal{N}(\ddot{q} | \mu_\theta(\mathbf{q}), \sigma^2_\theta(\mathbf{q})) $$

    \begin{block}{Training Objective: Negative Log-Likelihood (NLL)}
        Instead of MSE, we minimize the NLL to learn both the prediction and the uncertainty simultaneously:
        $$ \mathcal{L}_{NLL} = \frac{1}{2} \log \sigma^2_\theta(\mathbf{q}) + \frac{(\ddot{q}_{gt} - \mu_\theta(\mathbf{q}))^2}{2 \sigma^2_\theta(\mathbf{q})} + C $$
    \end{block}
    
    \begin{block}{Intuition}
        The model can reduce the penalty of a high error $(\ddot{x} - \mu)^2$ by increasing its uncertainty $\sigma^2$, effectively "admitting" it doesn't know the answer.
    \end{block}
    
    
\end{frame}

\begin{frame}{Neuro Probabilistic Model: Evolution}

    \begin{block}{Monte Carlo simulation}
        It's possible to use a \textbf{Monte Carlo} simulation to predict the confidence region during evolution. Given a \textbf{set of pendulums} at every time step the acceleration is sampled according to the model output distribution $\mu$ and $\sigma$.
    \end{block}

    \begin{figure}[h]
            \centering
            \includegraphics[width=0.6\linewidth]{09_Neuro_Prob/neuro_no_dropout.png} \\
            \vspace{-1em}
    \end{figure} 
    
        
\end{frame}


\begin{frame}{Neuro Probabilistic Model: Evolution with dropout}
    
    \begin{block}{Use of dropout}
        It is possible to use \textbf{dropout} to differentiate the model, so to have an ensamble of models. Note that dropout must remain active also during testing. 
    \end{block}

    \centering
    \begin{figure}[h]
            \centering
            \includegraphics[width=0.6\linewidth]{09_Neuro_Prob/neuro_dropout.png} \\
            \vspace{-1em}
    \end{figure}    
             
\end{frame}

